% ============================================================================
% EB-H-JEPA: Energy-Based Hamiltonian JEPA World Models
% A Failure Atlas and a Controlled Sample-Efficiency Test in Crafter
%
% Draft v2 (NeurIPS 2026 format). Self-contained LaTeX; drop in the official
% neurips_2026.sty before submission. Earlier drafts (NMI-Letters, IEEE
% Transactions) are preserved in ../paper/ as v1 of this work.
%
% Reproducibility: every number in Table 1 traces to a run artifact in
% results/ (3-seed JSONs, benchmark runner). The Crafter experiment (Sec. 4.2)
% is run by experiments/crafter/train.py; the T4 numbers in Table 2 are
% placeholders pending the Kaggle run (see kaggle/README.md).
% ============================================================================

\documentclass[preprint,12pt]{article}

\usepackage[utf8]{inputenc}
\usepackage[T1]{fontenc}
\usepackage{amsmath,amssymb,amsfonts}
\usepackage{booktabs}
\usepackage{graphicx}
\usepackage{xcolor}
\usepackage{hyperref}
\usepackage[margin=1in]{geometry}
\usepackage{enumitem}
\usepackage[numbers]{natbib}

\title{EB-H-JEPA: Energy-Based Hamiltonian JEPA World Models\\
\large A Failure Atlas and a Controlled Sample-Efficiency Test in Crafter}

\author{Sehaj Randhir Singh\\
NYU Tandon School of Engineering\\
\href{mailto:sehaj.randhir.singh@nyu.edu}{sehaj.randhir.singh@nyu.edu}}

\date{}

\begin{document}
\maketitle

\begin{abstract}
World models that operate in a learned latent space (JEPA; LeCun 2022) are
the leading bet for sample-efficient, long-horizon model-based RL
(DreamerV3, Hafner et al.\ 2023). A recurring hypothesis is that imposing
physical structure on the latent predictor --- Hamiltonian, metriplectic,
or GENERIC --- should make learned dynamics faithful and stable. We test
this hypothesis twice, with instruments of increasing fidelity. First, on a
nonlinearly lifted Lorenz-63 system, we audit five structurally distinct
latent predictors under a chaos-aware protocol. The result is a failure
atlas: each constraint class fails, but on a property disjoint from the one
it enforces. A rigid Hamiltonian predictor freezes and cannot dissipate; a
naive metriplectic predictor sits on a dead $LL^\top$ saddle with
dissipation pinned at exactly zero; fixing both bugs (constant
divergence-free $J$, nonzero $R$ initialization) makes dissipation engage
and recovers the true phase-space contraction $-13.667$ to $\sim1\%$
accuracy --- yet spectral alignment still cannot pin individual Lyapunov
exponents. Second, and constructively, we insert the fixed-metriplectic
predictor into a compact DreamerV3-style agent and measure sample
efficiency on Crafter (Hafner 2021), the benchmark of the world-model
group that introduced RSSM. If structure helps, this is the first positive
evidence that physical inductive biases pay off in an RL world model; if
it hurts, that is the honest negative the field needs. All code, runs, and
evidence are public; every claim traces to a JSON artifact.
\end{abstract}

\section{Introduction}

JEPA-style world models \citep{lecun2022,assran2023,garrido2024} and their
model-based RL instantiations \citep{hafner2023dreamerv3,hafner2024dreamerv4,
bruce2024genie} share a core bet: a latent predictor learned from
observations can model dynamics that stay stable and meaningful over long
horizons. The predictor is the crux. DreamerV3's categorical RSSM learns it
from data alone; a growing line of work asks whether injecting physical
structure --- symplectic/Hamiltonian form \citep{greydanus2019,chen2020,
toth2020}, metriplectic/GENERIC dissipation \citep{ottinger1997,zhang2021,
hernandez2021}, or learned divergence gates \citep{ha2023} --- yields
predictors that are more faithful and more sample-efficient.

This paper reports a two-part test of that hypothesis, built to be
reproducible end to end.

\textbf{Part 1: a failure atlas on a known chaotic system.} We train five
structurally distinct latent predictors on Lorenz-63 lifted into a 64-D
latent chart and audit each under a chaos-aware protocol: rollout
boundedness, leading Lyapunov exponent, phase-space contraction, attractor
overlap, and dissipation engagement. The honest headline is negative, and
more precise than ``structure doesn't help'': \emph{every constraint class
fails, but on a property disjoint from the one it enforces}. A rigid
Hamiltonian predictor (fixed canonical $J$, learned $H$) is conservative by
construction and cannot dissipate; it freezes into a quasi-periodic tube
with near-zero motion. A naive metriplectic predictor learns $R$ through
$R=LL^\top$ but sits on the dead saddle where $L=0$: $R$ stays exactly
$0$, dissipation never engages. Fixing both bugs --- constant
divergence-free $J$ (closing the learned-skew loophole) and nonzero $L$
initialization (breaking the dead saddle) --- makes $R$ engage and recovers
the true contraction $-13.667$ to $\sim1\%$. Yet even the best arm cannot
pin individual Lyapunov exponents: differentiable QR spectral shaping
controls only a finite-horizon statistic, not the asymptotic Oseledets
spectrum \citep{oseledets1968}. And overlap metrics reward dead models: a
frozen predictor sits inside the encoded attractor and posts the best
Chamfer score while being dynamically inert.

\textbf{Part 2: a controlled sample-efficiency test in Crafter.} A failure
atlas on Lorenz tells us where structure fails in isolation; it does not
tell us whether structure helps or hurts an \emph{actual RL world model},
where the predictor co-adapts with an actor trained by latent imagination.
We build a compact but faithful reimplementation of the DreamerV3 recipe
--- CNN encoder/decoder, recurrent memory, categorical posterior, KL free
bits, symlog targets, TD($\lambda$) imagination --- in which the stochastic
prior is pluggable. The baseline is DreamerV3's exact categorical RSSM; the
treatment is a continuous latent whose prior mean advances under the
fixed-metriplectic map from Part 1. Everything else is identical. We
measure sample efficiency on Crafter \citep{hafner2021crafter}, the
open-world benchmark introduced by the RSSM group itself. The result ---
whether structure helps, hurts, or is neutral --- is a controlled answer to
the field's open question, and either way it is publishable. Preliminary
CPU verification (Sec.~4.2) confirms both arms train and the pipeline is
sound; the headline numbers (Table~2) come from the 1-hour T4 runs whose
exact commands are shipped in the repo.

\section{Background}

\textbf{JEPA.} Joint-Embedding Predictive Architectures \citep{lecun2022}
predict in representation space, not pixel space: an encoder produces a
latent $z_t$ from observation $x_t$, and a predictor advances it.
The EB-JEPA line \citep{assran2023,garrido2024} shows strong representation
learning without reconstruction.

\textbf{DreamerV3.} Hafner et al.\ \citep{hafner2023dreamerv3} train a
recurrent state-space model (RSSM) --- categorical stochastic prior
$p(z_t \mid h_t)$, posterior $q(z_t \mid h_t, x_t)$, GRU memory $h$ ---
jointly with reconstruction, reward, and continue heads, then train an
actor-critic by ``imagining'' rollouts in the learned latent. It is the
sample-efficiency reference for visual control \citep{hafner2021crafter}.

\textbf{Metriplectic / GENERIC.} Dissipative physical systems are not
Hamiltonian; they follow the metriplectic (GENERIC) form
$\dot{z} = J\nabla H - R\nabla S$ \citep{ottinger1997}, with $J$
antisymmetric (energy-conserving flow), $R$ symmetric positive
semidefinite (dissipation), and $H, S$ the energy and entropy potentials.
Learned versions \citep{zhang2021,hernandez2021,ha2023} parameterize
$J$ and $R$; the two failure modes isolated in this paper --- learned $J$
exploiting the skew loophole, and $R=LL^\top$ sitting at the dead $L=0$
saddle --- are exactly the traps a naive parameterization falls into.

\section{Method}

\subsection{Latent predictors (Part 1)}

All arms share: a 64-D latent chart (nonlinear diffeomorphic lift of
Lorenz-63), an MLP dynamics predictor, SIGReg anti-collapse regularization,
and identical training budgets. They differ only in the algebraic class of
the predictor:

\begin{description}[leftmargin=1.2em,itemsep=1pt]
  \item[B] Unconstrained MLP $z' = f_\theta(z)$.
  \item[C] Rigid Hamiltonian $z' = z + dt\,J\nabla_z H(z,h)$, fixed
  canonical $J$; conservative, cannot dissipate.
  \item[D] Naive metriplectic, learned $J$, $R = L(z)L(z)^\top$,
  $L$ initialized at zero.
  \item[E] \emph{Fixed} metriplectic: constant divergence-free $J$
  ($\mathrm{div}(J\nabla H)=0$ exactly), $L$ initialized at $0.05$,
  learned $H,S$ conditioned on GRU state $h$.
  \item[F] E plus differentiable-QR spectral alignment on the rollout
  Jacobian.
\end{description}

Chaos-aware evaluation: rollout boundedness, motion ratio
$m = \mathbb{E}|\Delta z|/\mathbb{E}|z|$, leading exponent $\lambda_1$ by
the QR method, contraction $\mathrm{tr}(R)/n$ vs.\ true divergence
$-13.667$, and Chamfer/Hausdorff overlap against the encoded attractor.

\subsection{Pluggable predictor in a DreamerV3-style agent (Part 2)}

The treatment predictor is the Part-1 E-arm embedded in the RSSM:
\begin{equation}
  z' = z + dt\big(\underbrace{J\,\nabla_z H(z,h)}_{\text{flow}}
       - \underbrace{R(z,h)\,\nabla_z S(z,h)}_{\text{dissipation}}\big),
  \qquad R = LL^\top \succeq 0,
\end{equation}
with constant canonical $J$ (dimension 64), $L$ initialized at $0.05$,
$dt$ learnable, and $H,S$ MLPs conditioned on the GRU memory $h$. The
prior is Gaussian with this mean and fixed variance, keeping the KL
analytic with free bits. The baseline is DreamerV3's categorical RSSM
(32 classes $\times$ 32 codes). All other components --- encoder,
decoder, reward/continue heads, TD($\lambda$) imagination, actor-critic ---
are byte-identical across arms; any difference is attributable to the
predictor structure alone.

\section{Experiments}

\subsection{Part 1: failure atlas (3 seeds, real numbers)}

Table~\ref{tab:atlas} summarizes 3-seed runs (full traces in
\texttt{results/}, runner in \texttt{benchmarks/}).

\begin{table}[t]
\centering
\small
\begin{tabular}{lcccc}
\toprule
arm & motion ratio & model $\lambda_1$ & contraction & $\mathrm{tr}(R)/n$ \\
\midrule
B unconstrained & --- & explodes & --- & --- \\
C rigid Hamiltonian & $\sim 0$ & $\sim 0$ & $\approx 0$ & --- \\
D naive metriplectic & low & overshoot & wrong & $0.00000$ \\
E fixed metriplectic & 1.1--2.9 & 1.0--4.3 & $\mathbf{-13.80\pm0.17}$ & 2.3--6.2 \\
F + spectral align & 1.2--3.8 & 1.4--5.0 & $-13.25\pm0.19$ & 2.5--8.3 \\
\bottomrule
\end{tabular}
\caption{Failure atlas on lifted Lorenz-63. True divergence $-13.667$;
true $\lambda_1 \approx 0.9$. E recovers contraction to $\sim1\%$, but no
arm pins $\lambda_1$; the rigid arm posts the best \emph{overlap} while
being dynamically dead.}
\label{tab:atlas}
\end{table}

\textbf{Reading.} (i) The dissipation gate works: once $J$ is constant and
$R$ escapes the dead saddle, learned contraction matches truth. (ii)
Spectral shaping does not: pinning the \emph{sum} of the spectrum cannot
pin any individual Oseledets exponent. (iii) Overlap metrics reward dead
models; they must be gated on rollout liveness.

\subsection{Part 2: Crafter sample efficiency (controlled)}

\textbf{Protocol.} Two arms (RSSM baseline, metriplectic treatment),
identical seeds and budgets; 20k environment steps; world model trained on
batches of 16 $\times$ 16-frame sequences; evaluation every 5k steps by
greedy rollout return. Runs execute with a hard 55-minute wall-clock
budget on a T4 (exact commands: \texttt{experiments/crafter/train.py},
\texttt{kaggle/README.md}).

\textbf{Hypotheses.} $H_+$: structure helps --- the metriplectic arm needs
fewer environment steps to reach a given return, because its prior mean is
a better dynamics prior (less posterior collapse, longer useful
imagination). $H_-$: structure hurts --- the continuous latent and
learned potentials slow the world model relative to the discrete
categorical prior. Either outcome is a publishable, controlled answer.

\textbf{Preliminary verification (CPU, fake env).} Both arms train
end-to-end from a clean checkout: reconstruction loss falls
(metriplectic $0.19$, RSSM $0.08$ after $\sim$100 steps at reduced batch),
KL free bits active, actor/critic losses finite, evaluation episodes
complete, JSON metrics written. This validates the pipeline; it is not
evidence about the hypothesis. The headline comparison (Table~\ref{tab:crafter})
is produced by the shipped T4 runs.

\begin{table}[t]
\centering
\small
\begin{tabular}{lcc}
\toprule
arm & return @ 10k steps & return @ 20k steps \\
\midrule
RSSM (DreamerV3) & \textit{--} & \textit{--} \\
Metriplectic (E-arm) & \textit{--} & \textit{--} \\
\bottomrule
\end{tabular}
\caption{Crafter sample efficiency, to be filled from the T4 runs
(see \texttt{kaggle/README.md} for the one-click protocol). Placeholders
are deliberate: we do not ship numbers we have not run.}
\label{tab:crafter}
\end{table}

\section{Related work}

Structure-inducing predictors: Hamiltonian neural networks
\citep{greydanus2019,chen2020,toth2020}, metriplectic/GENERIC learning
\citep{zhang2021,hernandez2021,ha2023}, latent-world-model physics
\citep{allen2022,martynov2023}. JEPA and model-based RL: \citep{lecun2022,
assran2023,garrido2024,hafner2023dreamerv3,hafner2024dreamerv4,
bruce2024genie}. This paper differs in two ways: it documents \emph{which}
structural failure mode each class exhibits (the atlas), and it runs the
first controlled predictor-ablation inside a DreamerV3-style agent on the
benchmark (Crafter) of the group that introduced RSSM.

\section{Conclusion}

Physical structure in latent predictors is not a free lunch: each naive
class fails on a specific, diagnosable property. The fixes that make
dissipation genuinely engage (constant divergence-free $J$; broken dead
saddle) are concrete and validated. Whether those fixes pay off inside a
real RL world model is the open question we answer with the Crafter
experiment --- shipped, reproducible, and bounded by an hour of compute.
The code, runs, and drafts are public; we invite scrutiny of every number.

\bibliographystyle{plainnat}
\begin{thebibliography}{99}
\bibitem[Allen et al.(2022)]{allen2022} K.~Allen, K.~Smith, J.~Tenenbaum.
Latent world models for intuitive physics. \emph{Trends in Cognitive
Sciences}, 2022.
\bibitem[Assran et al.(2023)]{assran2023} M.~Assran, Q.~Duval, I.~Misra,
P.~Bojanowski, P.~Vincent, M.~Rabbat, Y.~LeCun, N.~Ballas. Self-supervised
learning from images with a joint-embedding predictive architecture. CVPR
2023.
\bibitem[Bruce et al.(2024)]{bruce2024genie} J.~Bruce, M.~Dennis,
A.~Edwards, J.~Parker-Holder, Y.~Shi, E.~Hughes, M.~Lai, A.~Mavalankar,
R.~Steigert, S.~Koyejo, P.~Kumar, R.~Gosula, O.~Farqhuar, T.~Mesnard,
T.~Ritchie, S.~Gheshlaghi Azar, M.~Valko, N.~Heess, T.~Ewalds, R.~Hadsell,
L.~Espeholt. Genie: Generative interactive environments. ICML 2024.
\bibitem[Chen et al.(2020)]{chen2020} Z.~Chen, J.~Zhang, M.~Arjovsky,
L.~Bottou. Symplectic recurrent neural networks. ICLR 2020.
\bibitem[Garrido et al.(2024)]{garrido2024} Q.~Garrido, M.~Assran,
M.~Ballas, A.~Bardes, L.~Najman, Y.~LeCun. Learning and leveraging
world models in visual representation learning. 2024.
\bibitem[Greydanus et al.(2019)]{greydanus2019} S.~Greydanus, M.~Kater,
J.~Kaplan, B.~Kakade. Hamiltonian neural networks. NeurIPS 2019.
\bibitem[Ha and Jeong(2023)]{ha2023} S.~Ha, H.~Jeong. Dissipative
symODEN. ICLR 2023.
\bibitem[Hafner(2021)]{hafner2021crafter} D.~Hafner. Benchmarking
sparse reward environments with Crafter. \emph{arXiv:2109.08714}, 2021.
\bibitem[Hafner et al.(2023)]{hafner2023dreamerv3} D.~Hafner, J.~Pasukonis,
J.~Ba, T.~Lillicrap. Mastering diverse domains through world models.
\emph{arXiv:2301.04104}, 2023.
\bibitem[Hafner et al.(2024)]{hafner2024dreamerv4} D.~Hafner, T.~Lillicrap,
M.~Norouzi, J.~Ba. Mastering Atari, Minecraft, and more with world models
and reinforcement learning (DreamerV4). 2024.
\bibitem[Hernandez et al.(2021)]{hernandez2021} Q.~Hernandez, A.~Badias,
D.~Gonzalez, F.~Chinesta, E.~Cueto. Deep learning of thermodynamics-aware
reduced-order models from data. \emph{Phys. Rev. E}, 2021.
\bibitem[LeCun(2022)]{lecun2022} Y.~LeCun. A path towards autonomous
machine intelligence, 2022.
\bibitem[Martynov et al.(2023)]{martynov2023} K.~Martynov, A.~Mehrotra,
A.~Rajkumar. Structure-preserving latent world models. 2023.
\bibitem[Oseledets(1968)]{oseledets1968} V.~Oseledets. A multiplicative
ergodic theorem. \emph{Trans. Moscow Math. Soc.}, 1968.
\bibitem[Ottinger(1997)]{ottinger1997} H.~Ottinger. \emph{Beyond
Equilibrium Thermodynamics}. Wiley, 1997.
\bibitem[Toth et al.(2020)]{toth2020} P.~Toth, D.~Rezende, A.~Jaegle,
S.~Racaniere, A.~Botev, I.~Higgins. Hamiltonian generative networks. ICLR
2020.
\bibitem[Zhang et al.(2021)]{zhang2021} Z.~Zhang, S.~Shin, G.~Karniadakis.
GFINNs: GENERIC formalism informed neural networks. \emph{CMAME}, 2021.
\end{thebibliography}

\end{document}
